\documentclass[10pt,a4paper]{amsart}
\usepackage{amsmath,amssymb,amsthm,graphicx}
% For commutative diagrams we use Xy-pic
\usepackage[all]{xy}

% We use 2cell for 2-commutative diagrams.
\xyoption{2cell}
% Theorem environments.
%
\theoremstyle{plain}
\newtheorem{theorem}[subsection]{Theorem}
\newtheorem{proposition}[subsection]{Proposition}
\newtheorem{lemma}[subsection]{Lemma}

\theoremstyle{definition}
\newtheorem{definition}[subsection]{Definition}
\newtheorem{example}[subsection]{Example}
\newtheorem{exercise}[subsection]{Exercise}
\newtheorem{situation}[subsection]{Situation}

\theoremstyle{remark}
\newtheorem{remark}[subsection]{Remark}
\newtheorem{remarks}[subsection]{Remarks}

\numberwithin{equation}{subsection}

% Macros
%
\def\lim{\mathop{\mathrm{lim}}\nolimits}
\def\colim{\mathop{\mathrm{colim}}\nolimits}
\def\Spec{\mathop{\mathrm{Spec}}}
\def\Hom{\mathop{\mathrm{Hom}}\nolimits}
\def\Ext{\mathop{\mathrm{Ext}}\nolimits}
\def\SheafHom{\mathop{\mathcal{H}\!\mathit{om}}\nolimits}
\def\SheafExt{\mathop{\mathcal{E}\!\mathit{xt}}\nolimits}
\def\Sch{\mathit{Sch}}
\def\Mor{\mathop{\mathrm{Mor}}\nolimits}
\def\Ob{\mathop{\mathrm{Ob}}\nolimits}
\def\Sh{\mathop{\mathit{Sh}}\nolimits}
\def\NL{\mathop{N\!L}\nolimits}
\def\CH{\mathop{\mathrm{CH}}\nolimits}
\def\proetale{{pro\text{-}\acute{e}tale}}
\def\etale{{\acute{e}tale}}
\def\QCoh{\mathit{QCoh}}
\def\Ker{\mathop{\mathrm{Ker}}}
\def\Im{\mathop{\mathrm{Im}}}
\def\Coker{\mathop{\mathrm{Coker}}}
\def\Coim{\mathop{\mathrm{Coim}}}

% Boxtimes
%
\DeclareMathSymbol{\boxtimes}{\mathbin}{AMSa}{"02}

%
% Macros for moduli stacks/spaces
%
\def\QCohstack{\mathcal{QC}\!\mathit{oh}}
\def\Cohstack{\mathcal{C}\!\mathit{oh}}
\def\Spacesstack{\mathcal{S}\!\mathit{paces}}
\def\Quotfunctor{\mathrm{Quot}}
\def\Hilbfunctor{\mathrm{Hilb}}
\def\Curvesstack{\mathcal{C}\!\mathit{urves}}
\def\Polarizedstack{\mathcal{P}\!\mathit{olarized}}
\def\Complexesstack{\mathcal{C}\!\mathit{omplexes}}
% \Pic is the operator that assigns to X its picard group, usage \Pic(X)
% \Picardstack_{X/B} denotes the Picard stack of X over B
% \Picardfunctor_{X/B} denotes the Picard functor of X over B
\def\Pic{\mathop{\mathrm{Pic}}\nolimits}
\def\Picardstack{\mathcal{P}\!\mathit{ic}}
\def\Picardfunctor{\mathrm{Pic}}
\def\Deformationcategory{\mathcal{D}\!\mathit{ef}}
\usepackage{cleveref}
\crefname{originallemma}{original lemma}{original lemmas}
\Crefname{originallemma}{Original lemma}{Original lemmas}
\makeatletter
\newlabel{limits-section-limits}{{01YV}{source}{}{sourcepointer-1}{}}
\newlabel{limits-section-limits@cref}{{[originallemma][0][]01YV}{[1][0][]source}{}{}{}}
\makeatother
\crefname{originallemma}{original lemma}{original lemmas}
\Crefname{originallemma}{Original lemma}{Original lemmas}
\makeatletter
\newlabel{derived-lemma-negative-vanishing}{{05TC}{source}{}{sourcepointer-2}{}}
\newlabel{derived-lemma-negative-vanishing@cref}{{[originallemma][0][]05TC}{[1][0][]source}{}{}{}}
\makeatother
\title{Higher direct images along a limit of morphisms, for a coefficient complex}
\author{The Stacks Project Authors}
\date{stacks-project:etale-cohomology.tex:local-snapshot-20261011; GNU FDL 1.2 or later}
\begin{document}
\maketitle
\noindent\textit{Editorial note.} All delivered mathematical text is the author's own, copied from the pinned Stacks Project chapter file etale-cohomology.tex; no statement or proof step was rewritten, reordered or removed.\par\medskip
\noindent\textit{Editorial note.} The proof is the author's complete argument: the reduction to cohomology sheaves and the reassembly are both inside the counted span.\par\medskip
\noindent\textit{Editorial note.} Every labelled reference occurring inside the counted spans is carried in assembly.external\_references as a non-load-bearing navigation pointer (same-chapter ones at their definition site, cross-chapter ones at the invocation inside the counted span), so no load-bearing mathematics is replaced by a stub. No bibliography citation occurs inside the counted spans.\par\medskip
\noindent\textit{Editorial note.} The checked chapter list for same-name collisions is the set of stacks chapters staged in this repository: chow, cohomology, divisors, etale-cohomology, flat, more-algebra, more-morphisms, morphisms, spaces-more-morphisms, varieties together with this card's chapter etale-cohomology.\par\medskip
\noindent\textit{Editorial note.} The preamble is the chapter's own native preamble excerpt already used for this source: the Xy-pic setup (native preamble lines 14-18) and the native theorem declarations with the mathematical macros (native preamble lines 154-217).\par\medskip
\noindent\textit{Allowed standard prerequisites (editorial).}\begin{itemize}
\item The set-up of the companion result: inverse system of morphisms with affine transitions and the limit morphism. Same chapter, the statement of lemma-relative-colimit-general delivered verbatim as a context span in this card (source lines 7566-7582). Supplies the notation and the hypotheses, and the sheaf-level statement which the complex-level statement generalises.
\item Spectral sequence or resolution argument comparing the cohomology sheaves of the pushforward. Same chapter, the technique by which a statement about sheaves is upgraded to bounded below complexes, together with the smoothness of the pullback functors. Supplies the reduction from the complex case to the sheaf case.
\end{itemize}\medskip
\ifdefined\hypertarget\hypertarget{sourcepointer-1}{}\fi
\noindent\textit{Source reference.} Original lemma 01YV: sources/stacks-etale-cohomology/pinned-original-source.tex, source lines 7573--7573. Navigation pointer invoked at proof line 7573; the counted span contains the pointer only and the cited result is not reproduced inside it, so the pointer is registered as a non-load-bearing navigation stub in the established form for Stacks material. No bibliography citation is involved.\par\medskip
\ifdefined\hypertarget\hypertarget{sourcepointer-2}{}\fi
\noindent\textit{Source reference.} Original lemma 05TC: sources/stacks-etale-cohomology/pinned-original-source.tex, source lines 7896--7896. Navigation pointer invoked at proof line 7896; the counted span contains the pointer only and the cited result is not reproduced inside it, so the pointer is registered as a non-load-bearing navigation stub in the established form for Stacks material. No bibliography citation is involved.\par\medskip
\begin{lemma}
\label{lemma-relative-colimit-general}
Let $I$ be a directed set. Let $g_i : X_i \to S_i$ be an inverse system of
morphisms of schemes over $I$. Assume $g_i$ is quasi-compact and
quasi-separated and for $i' \geq i$ the transition morphisms
$f_{i'i} : X_{i'} \to X_i$ and $h_{i'i} : S_{i'} \to S_i$ are affine.
Let $g : X \to S$ be the limit of the morphisms $g_i$, see
Limits, Section \ref{limits-section-limits}.
Denote $f_i : X \to X_i$ and $h_i : S \to S_i$ the projections.
Let $(\mathcal{F}_i, \varphi_{i'i})$ be a system of sheaves
on $(X_i, f_{i'i})$. Set $\mathcal{F} = \colim f_i^{-1}\mathcal{F}_i$. Then
$$
R^p g_* \mathcal{F} =
\colim_{i \in I} h_i^{-1}R^p g_{i, *} \mathcal{F}_i
$$
for all $p \geq 0$.
\end{lemma}

\begin{lemma}
\label{lemma-relative-colimit-general-complexes}
Let $I$, $g_i : X_i \to S_i$, $g : X \to S$, $f_i$, $g_i$, $h_i$ be as in
Lemma \ref{lemma-relative-colimit-general}.
Let $0 \in I$ and $K_0 \in D^+(X_{0, \etale})$.
For $i \geq 0$ denote $K_i$ the pullback of $K_0$ to $X_i$.
Denote $K$ the pullback of $K_0$ to $X$. Then
$$
R^pg_*K = \colim_{i \geq 0} h_i^{-1}R^pg_{i, *}K_i
$$
for all $p \in \mathbf{Z}$.
\end{lemma}

\begin{proof}
Fix an integer $p_0 \in \mathbf{Z}$.
Let $a$ be an integer such that $H^j(K_0) = 0$ for $j < a$.
We will prove the formula holds for all $p \leq p_0$
by descending induction on $a$. If
$a > p_0$, then we see that the left and right hand side of
the formula are zero for $p \leq p_0$ by trivial vanishing, see
Derived Categories, Lemma \ref{derived-lemma-negative-vanishing}.
Assume $a \leq p_0$. Consider the distinguished triangle
$$
H^a(K_0)[-a] \to K_0 \to \tau_{\geq a + 1}K_0
$$
Pulling back this distinguished triangle to $X_i$ and $X$
gives compatible distinguished triangles for $K_i$ and $K$.
For $p \leq p_0$ we consider the commutative diagram
$$
\xymatrix{
\colim_{i \geq 0} h_i^{-1}R^{p - 1}g_{i, *}(\tau_{\geq a + 1}K_i)
\ar[r]_-\alpha \ar[d] &
R^{p - 1}g_*(\tau_{\geq a + 1}K) \ar[d] \\
\colim_{i \geq 0} h_i^{-1}R^pg_{i, *}(H^a(K_i)[-a])
\ar[r]_-\beta \ar[d] &
R^pg_*(H^a(K)[-a]) \ar[d] \\
\colim_{i \geq 0} h_i^{-1}R^pg_{i, *}K_i
\ar[r]_-\gamma \ar[d] &
R^pg_*K \ar[d] \\
\colim_{i \geq 0} R^pg_{i, *}\tau_{\geq a + 1}K_i
\ar[r]_-\delta \ar[d] &
R^pg_*\tau_{\geq a + 1}K \ar[d] \\
\colim_{i \geq 0} R^{p + 1}g_{i, *}(H^a(K_i)[-a])
\ar[r]^-\epsilon &
R^{p + 1}g_*(H^a(K)[-a])
}
$$
with exact columns. The arrows $\beta$ and $\epsilon$ are isomorphisms by
Lemma \ref{lemma-relative-colimit-general}.
The arrows $\alpha$ and $\delta$ are isomorphisms
by induction hypothesis.
Hence $\gamma$ is an isomorphism as desired.
\end{proof}

\end{document}
